\documentclass[12pt]{article}

% Idioma y codificación
\usepackage[spanish, es-tabla, es-notilde]{babel}       %es-tabla para que se titule "Tabla"
\usepackage[utf8]{inputenc}

% Márgenes
\usepackage[a4paper,top=3cm,bottom=2.5cm,left=3cm,right=3cm]{geometry}

% Comentarios de bloque
\usepackage{verbatim}

% Paquetes de links
\usepackage[hidelinks]{hyperref}    % Permite enlaces
\usepackage{url}                    % redirecciona a la web

% Más opciones para enumeraciones
\usepackage{enumitem}

% Personalizar la portada
\usepackage{titling}

% Paquetes de tablas
\usepackage{multirow}

% Para añadir el símbolo de euro
\usepackage{eurosym}

% Paquete para subrayar con diferentes estilos
\usepackage[normalem]{ulem}


%------------------------------------------------------------------------

%Paquetes de figuras
\usepackage{caption}
\usepackage{subcaption} % Figuras al lado de otras
\usepackage{float}      % Poner figuras en el sitio indicado H.


% Paquetes de imágenes
\usepackage{graphicx}       % Paquete para añadir imágenes
\usepackage{transparent}    % Para manejar la opacidad de las figuras

% Paquete para usar colores
\usepackage[dvipsnames, table, xcdraw]{xcolor}
\usepackage{pagecolor}      % Para cambiar el color de la página

% Habilita tamaños de fuente mayores
\usepackage{fix-cm}

% Para los gráficos
\usepackage{tikz}
\usepackage{forest}

% Para poder situar los nodos en los grafos
\usetikzlibrary{positioning}

% Para evitar problemas entre babel y tikz
\usetikzlibrary{babel}


%------------------------------------------------------------------------

% Paquetes de matemáticas
\usepackage{mathtools, amsfonts, amssymb, mathrsfs}
\usepackage[makeroom]{cancel}     % Simplificar tachando
\usepackage{polynom}    % Divisiones y Ruffini
\usepackage{units} % Para poner fracciones diagonales con \nicefrac

\usepackage{pgfplots}   %Representar funciones
\pgfplotsset{compat=1.18}  % Versión 1.18

\usepackage{tikz-cd}    % Para usar diagramas de composiciones
\usetikzlibrary{calc}   % Para usar cálculo de coordenadas en tikz

%Definición de teoremas, etc.
\usepackage{amsthm}
%\swapnumbers   % Intercambia la posición del texto y de la numeración

\theoremstyle{plain}

\makeatletter
\@ifclassloaded{article}{
  % --- Configuración para ARTICLE ---
  \newtheorem{teo}{Teorema}[section]
  \theoremstyle{definition}
  \newtheorem{definicion}{Definición}[section]
  \newtheorem{ejercicio}{Ejercicio}[section]

  \renewcommand{\theteo}{%
    \ifnum\value{section}=0 \arabic{teo}\else \thesection.\arabic{teo}\fi
  }
  \renewcommand{\thedefinicion}{%
    \ifnum\value{section}=0 \arabic{definicion}\else \thesection.\arabic{definicion}\fi
  }
  \renewcommand{\theejercicio}{%
    \ifnum\value{section}=0 \arabic{ejercicio}\else \thesection.\arabic{ejercicio}\fi
  }
}{
  % --- Configuración para BOOK / REPORT ---
  \newtheorem{teo}{Teorema}[chapter]
  \theoremstyle{definition}
  \newtheorem{definicion}{Definición}[chapter]
  \newtheorem{ejercicio}{Ejercicio}[section]

  % Si no hay sección: Capítulo.Teorema (ej. 2.1)
  % Si hay sección:    Capítulo.Sección.Teorema (ej. 2.3.1)
  \renewcommand{\theteo}{%
    \ifnum\value{section}=0
      \thechapter.\arabic{teo}%
    \else
      \thesection.\arabic{teo}% (\thesection ya contiene el número de capítulo)
    \fi
  }
  \renewcommand{\thedefinicion}{%
    \ifnum\value{section}=0
      \thechapter.\arabic{definicion}%
    \else
      \thesection.\arabic{definicion}%
    \fi
  }
  \renewcommand{\theejercicio}{%
    \ifnum\value{section}=0
      \arabic{ejercicio}%
    \else
      \thesection.\arabic{ejercicio}% Formato por defecto: ej. 2.2.1
    \fi
  }
}
\makeatother

% Teoremas dependientes o con contador compartido
\newtheorem{coro}{Corolario}[teo]
\newtheorem{prop}[teo]{Proposición}
\newtheorem{lema}[teo]{Lema}
\newtheorem*{lema*}{Lema}

\theoremstyle{remark}
\newtheorem*{observacion}{Observación}

\theoremstyle{definition}
\newtheorem*{notacion}{Notación}
\newtheorem*{ejemplo}{Ejemplo}
\newtheorem*{ejercicio*}{Ejercicio}


% \renewcommand\qedsymbol{$\blacksquare$}         % Cambiar símbolo QED
%------------------------------------------------------------------------

% Paquetes para encabezados
\usepackage{fancyhdr}
\pagestyle{fancy}
\fancyhf{}

\newcommand{\helv}{ % Modificación tamaño de letra
\fontfamily{}\fontsize{12}{12}\selectfont}
\setlength{\headheight}{15pt} % Amplía el tamaño del índice


%\usepackage{lastpage}   % Referenciar última pag   \pageref{LastPage}
%\fancyfoot[C]{%
%  \begin{minipage}{\textwidth}
%    \centering
%    ~\\
%    \thepage\\
%    \href{https://losdeldgiim.github.io/}{\texttt{\footnotesize losdeldgiim.github.io}}
%  \end{minipage}
%}
\fancyfoot[C]{\thepage}
\fancyfoot[R]{\href{https://losdeldgiim.github.io/}{\texttt{\footnotesize losdeldgiim.github.io}}}

%------------------------------------------------------------------------

% Conseguir que no ponga "Capítulo 1". Sino solo "1."
\makeatletter
\@ifclassloaded{book}{
  \renewcommand{\chaptermark}[1]{\markboth{\thechapter.\ #1}{}} % En el encabezado
    
  \renewcommand{\@makechapterhead}[1]{%
  \vspace*{50\p@}%
  {\parindent \z@ \raggedright \normalfont
    \ifnum \c@secnumdepth >\m@ne
      \huge\bfseries \thechapter.\hspace{1em}\ignorespaces
    \fi
    \interlinepenalty\@M
    \Huge \bfseries #1\par\nobreak
    \vskip 40\p@
  }}
}
\makeatother

%------------------------------------------------------------------------
% Paquetes de cógido
\usepackage{minted}
\renewcommand\listingscaption{Código fuente}

\usepackage{fancyvrb}
% Personaliza el tamaño de los números de línea
\renewcommand{\theFancyVerbLine}{\small\arabic{FancyVerbLine}}

% Estilo para C++
\newminted{cpp}{
    frame=lines,
    framesep=2mm,
    baselinestretch=1.2,
    linenos,
    escapeinside=||
}

% para minted
\definecolor{LightGray}{rgb}{0.95,0.95,0.92}
\setminted{
    linenos=true,
    stepnumber=5,
    numberfirstline=true,
    autogobble,
    breaklines=true,
    breakautoindent=true,
    breaksymbolleft=,
    breaksymbolright=,
    breaksymbolindentleft=0pt,
    breaksymbolindentright=0pt,
    breaksymbolsepleft=0pt,
    breaksymbolsepright=0pt,
    fontsize=\footnotesize,
    bgcolor=LightGray,
    numbersep=10pt
}


\usepackage{listings} % Para incluir código desde un archivo

\renewcommand\lstlistingname{Código Fuente}
\renewcommand\lstlistlistingname{Índice de Códigos Fuente}

% Definir colores
\definecolor{vscodepurple}{rgb}{0.5,0,0.5}
\definecolor{vscodeblue}{rgb}{0,0,0.8}
\definecolor{vscodegreen}{rgb}{0,0.5,0}
\definecolor{vscodegray}{rgb}{0.5,0.5,0.5}
\definecolor{vscodebackground}{rgb}{0.97,0.97,0.97}
\definecolor{vscodelightgray}{rgb}{0.9,0.9,0.9}

% Configuración para el estilo de C similar a VSCode
\lstdefinestyle{vscode_C}{
  backgroundcolor=\color{vscodebackground},
  commentstyle=\color{vscodegreen},
  keywordstyle=\color{vscodeblue},
  numberstyle=\tiny\color{vscodegray},
  stringstyle=\color{vscodepurple},
  basicstyle=\scriptsize\ttfamily,
  breakatwhitespace=false,
  breaklines=true,
  captionpos=b,
  keepspaces=true,
  numbers=left,
  numbersep=5pt,
  showspaces=false,
  showstringspaces=false,
  showtabs=false,
  tabsize=2,
  frame=tb,
  framerule=0pt,
  aboveskip=10pt,
  belowskip=10pt,
  xleftmargin=10pt,
  xrightmargin=10pt,
  framexleftmargin=10pt,
  framexrightmargin=10pt,
  framesep=0pt,
  rulecolor=\color{vscodelightgray},
  backgroundcolor=\color{vscodebackground},
}

%------------------------------------------------------------------------
% La etiqueda IA sirve para indicar a partir de donde le pasamos el preambulo a la IA para el escaneo
% IA

% Comandos definidos
\newcommand{\bb}[1]{\mathbb{#1}}
\newcommand{\cc}[1]{\mathcal{#1}}

% I prefer the slanted \leq
\let\oldleq\leq % save them in case they're every wanted
\let\oldgeq\geq
\renewcommand{\leq}{\leqslant}
\renewcommand{\geq}{\geqslant}

% Si y solo si
\newcommand{\sii}{\iff}

% MCD y MCM
\DeclareMathOperator{\mcd}{mcd}
\DeclareMathOperator{\mcm}{mcm}

% Signo
\DeclareMathOperator{\sgn}{sgn}

% Imagen
\DeclareMathOperator{\Ima}{Im}

% Aplicación identidad
\newcommand{\id}{\mathrm{Id}}

% Letras griegas
\newcommand{\eps}{\epsilon}
\newcommand{\veps}{\varepsilon}
\newcommand{\lm}{\lambda}

\newcommand{\ol}{\overline}
\newcommand{\ul}{\underline}
\newcommand{\wt}{\widetilde}
\newcommand{\wh}{\widehat}

\let\oldvec\vec
\renewcommand{\vec}{\overrightarrow}

% Derivadas parciales
\newcommand{\del}[2]{\frac{\partial #1}{\partial #2}}
\newcommand{\Del}[3]{\frac{\partial^{#1} #2}{\partial #3^{#1}}}
\newcommand{\deld}[2]{\dfrac{\partial #1}{\partial #2}}
\newcommand{\Deld}[3]{\dfrac{\partial^{#1} #2}{\partial #3^{#1}}}


\newcommand{\AstIg}{\stackrel{(\ast)}{=}}
\newcommand{\Hop}{\stackrel{L'H\hat{o}pital}{=}}

\newcommand{\red}[1]{{\color{red}#1}} % Para integrales, destacar los cambios.

% Método de integración
\newcommand{\MetInt}[2]{
    \left[\begin{array}{c}
        #1 \\ #2
    \end{array}\right]
}

% Declarar aplicaciones
% 1. Nombre aplicación
% 2. Dominio
% 3. Codominio
% 4. Variable
% 5. Imagen de la variable
\newcommand{\Func}[5]{
    \begin{equation*}
        \begin{array}{rrll}
            \displaystyle #1:& \displaystyle  #2 & \longrightarrow & \displaystyle  #3\\
               & \displaystyle  #4 & \longmapsto & \displaystyle  #5
        \end{array}
    \end{equation*}
}

%------------------------------------------------------------------------


\newcommand{\myhref}[2]{\href{#1}{\color{blue}\uline{#2}}}


\begin{document}

    % 1. Foto de fondo
    % 2. Título
    % 3. Encabezado Izquierdo
    % 4. Color de fondo
    % 5. Coord x del titulo
    % 6. Coord y del titulo
    % 7. Fecha

    
    % 1. Foto de fondo
% 2. Título
% 3. Encabezado Izquierdo
% 4. Color de fondo
% 5. Coord x del titulo
% 6. Coord y del titulo
% 7. Fecha
% 8. Autor

\newcommand{\portada}[8]{

    \portadaBase{#1}{#2}{#3}{#4}{#5}{#6}{#7}{#8}
    \portadaBook{#1}{#2}{#3}{#4}{#5}{#6}{#7}{#8}
}

\newcommand{\portadaExamen}[8]{

    \portadaBase{#1}{#2}{#3}{#4}{#5}{#6}{#7}{#8}
    \portadaArticle{#1}{#2}{#3}{#4}{#5}{#6}{#7}{#8}
}




\newcommand{\portadaBase}[8]{

    % Tiene la portada principal y la licencia Creative Commons
    
    % 1. Foto de fondo
    % 2. Título
    % 3. Encabezado Izquierdo
    % 4. Color de fondo
    % 5. Coord x del titulo
    % 6. Coord y del titulo
    % 7. Fecha
    % 8. Autor
    
    
    \thispagestyle{empty}               % Sin encabezado ni pie de página
    \newgeometry{margin=0cm}        % Márgenes nulos para la primera página
    
    
    % Encabezado
    \fancyhead[L]{\helv #3}
    \fancyhead[R]{\helv \nouppercase{\leftmark}}
    
    
    \pagecolor{#4}        % Color de fondo para la portada
    
    \begin{figure}[p]
        \centering
        \transparent{0.3}           % Opacidad del 30% para la imagen
        
        \includegraphics[width=\paperwidth, keepaspectratio]{../../_assets/#1}
    
        \begin{tikzpicture}[remember picture, overlay]
            \node[anchor=north west, text=white, opacity=1, font=\fontsize{60}{90}\selectfont\bfseries\sffamily, align=left] at (#5, #6) {#2};
            
            \node[anchor=south east, text=white, opacity=1, font=\fontsize{12}{18}\selectfont\sffamily, align=right] at (9.7, 3) {\href{https://losdeldgiim.github.io/}{\textbf{Los Del DGIIM}, \texttt{\footnotesize losdeldgiim.github.io}}};
            
            \node[anchor=south east, text=white, opacity=1, font=\fontsize{12}{15}\selectfont\sffamily, align=right] at (9.7, 1.8) {Doble Grado en Ingeniería Informática y Matemáticas\\Universidad de Granada};
        \end{tikzpicture}
    \end{figure}
    
    
    \restoregeometry        % Restaurar márgenes normales para las páginas subsiguientes
    \pagecolor{white}       % Restaurar el color de página
    
    
    \newpage
    \thispagestyle{empty}               % Sin encabezado ni pie de página
    \begin{tikzpicture}[remember picture, overlay]
        \node[anchor=south west, inner sep=3cm] at (current page.south west) {
            \begin{minipage}{0.5\paperwidth}
                \href{https://creativecommons.org/licenses/by-nc-nd/4.0/}{
                    \includegraphics[height=2cm]{../../_assets/Licencia.png}
                }\vspace{1cm}\\
                Esta obra está bajo una
                \href{https://creativecommons.org/licenses/by-nc-nd/4.0/}{
                    Licencia Creative Commons Atribución-NoComercial-SinDerivadas 4.0 Internacional (CC BY-NC-ND 4.0).
                }\\
    
                Eres libre de compartir y redistribuir el contenido de esta obra en cualquier medio o formato, siempre y cuando des el crédito adecuado a los autores originales y no persigas fines comerciales. 
            \end{minipage}
        };
    \end{tikzpicture}
    
    
    
    % 1. Foto de fondo
    % 2. Título
    % 3. Encabezado Izquierdo
    % 4. Color de fondo
    % 5. Coord x del titulo
    % 6. Coord y del titulo
    % 7. Fecha
    % 8. Autor


}


\newcommand{\portadaBook}[8]{

    % 1. Foto de fondo
    % 2. Título
    % 3. Encabezado Izquierdo
    % 4. Color de fondo
    % 5. Coord x del titulo
    % 6. Coord y del titulo
    % 7. Fecha
    % 8. Autor

    % Personaliza el formato del título
    \pretitle{\begin{center}\bfseries\fontsize{42}{56}\selectfont}
    \posttitle{\par\end{center}\vspace{2em}}
    
    % Personaliza el formato del autor
    \preauthor{\begin{center}\Large}
    \postauthor{\par\end{center}\vfill}
    
    % Personaliza el formato de la fecha
    \predate{\begin{center}\huge}
    \postdate{\par\end{center}\vspace{2em}}
    
    \title{#2}
    \author{\href{https://losdeldgiim.github.io/}{Los Del DGIIM, \texttt{\large losdeldgiim.github.io}}
    \\ \vspace{0.5cm}#8}
    \date{Granada, #7}
    \maketitle
    
    \tableofcontents
}




\newcommand{\portadaArticle}[8]{

    % 1. Foto de fondo
    % 2. Título
    % 3. Encabezado Izquierdo
    % 4. Color de fondo
    % 5. Coord x del titulo
    % 6. Coord y del titulo
    % 7. Fecha
    % 8. Autor

    % Personaliza el formato del título
    \pretitle{\begin{center}\bfseries\fontsize{42}{56}\selectfont}
    \posttitle{\par\end{center}\vspace{2em}}
    
    % Personaliza el formato del autor
    \preauthor{\begin{center}\Large}
    \postauthor{\par\end{center}\vspace{3em}}
    
    % Personaliza el formato de la fecha
    \predate{\begin{center}\huge}
    \postdate{\par\end{center}\vspace{5em}}
    
    \title{#2}
    \author{\href{https://losdeldgiim.github.io/}{Los Del DGIIM, \texttt{\large losdeldgiim.github.io}}
    \\ \vspace{0.5cm}#8}
    \date{Granada, #7}
    \thispagestyle{empty}               % Sin encabezado ni pie de página
    \maketitle
    \vfill
}
    \portadaExamen{ffccA4.jpg}{Análisis Funcional\\Examen XVII}{Análisis Funcional. Examen XVII}{MidnightBlue}{-8}{28}{2026}{Arturo Olivares Martos}

    \begin{description}
        \item[Asignatura] Análisis Funcional.
        \item[Curso Académico] 2024/2025.
        \item[Grado] Matemáticas.
        % \item[Grupo] ---.
        \item[Profesor] Miguel Martín Suárez.
        \item[Descripción] Parcial 2.
        \item[Fecha] 4 de diciembre de 2024.
        % \item[Duración] ---.
    
    \end{description}
    \newpage


    % ------------------------------------
    
    \begin{ejercicio}
        En el espacio de Banach $X = (C\left[0, 1\right], \|\cdot\|_\infty)$ se considera el funcional lineal $\varphi$ definido por
        \[
            \varphi(f) = \int_0^1 (2t - 1) f(t) \, dt \quad \forall f \in C\left[0, 1\right].
        \]
        \begin{enumerate}[label=\alph*)]
            \item Pruebe que $\varphi$ es continuo y calcule su norma.
            \item ¿Alcanza $\varphi$ su norma? Justifique su respuesta.
            \item ¿Existe $x^{**} \in X^{**}$ con $\|x^{**}\| = 1$ tal que $x^{**}(\varphi) = \|\varphi\|$? Justifique su respuesta.
        \end{enumerate}
    \end{ejercicio}

    \begin{ejercicio}
        Rodee las afirmaciones verdaderas (que pueden ser una, dos, tres o cuatro):
        \begin{enumerate}[label=\alph*)]
            \item Sea $X$ un subespacio cerrado de $\ell_{3000}$ de dimensión infinita ¿Puede ser $X$ isomorfo a $\ell_1$?
            \begin{enumerate}[label={a\arabic*)}]
                \item Si, porque $\ell_1 \subseteq \ell_{3000}$
                \item No, porque $\ell_1 \nsubseteq \ell_{3000}$
                \item No, porque $\ell_1$ es separable y $X$ no lo es
                \item No, porque $X^*$ es separable y $\ell_1^*$ no lo es
            \end{enumerate}
            \item Sea $\psi \colon \ell_2 \to \bb{R}$ el funcional lineal dado por $\psi(x) = \sum\limits_{n=1}^\infty \frac{x(n)}{n}$ para $x \in \ell_2$:
            \begin{enumerate}[label={b\arabic*)}]
                \item $\psi$ alcanza su norma porque $\ell_2$ es reflexivo
                \item $\psi$ alcanza su norma porque $\ell_2$ es un espacio de Hilbert
                \item $\psi$ no alcanza su norma
                \item $\psi$ no alcanza su norma porque $\ell_2$ no es completo
            \end{enumerate}
            \item Sea $Y = \bb{K}^{(\bb{N})}$ (el espacio vectorial de las sucesiones casi nulas):
            \begin{enumerate}[label={c\arabic*)}]
                \item la norma de la suma y la norma del máximo son equivalentes en $Y$
                \item la norma de la suma y la norma del máximo no son equivalentes en $Y$
                \item cualesquiera dos normas definidas en $Y$ son equivalentes
                \item no se puede definir ninguna norma en $Y$
            \end{enumerate}
            \item En un espacio normado $X$ tomamos un conjunto convexo y cerrado $C$ y un punto $x_0 \notin C$:
            \begin{enumerate}[label={d\arabic*)}]
                \item existe $f \in X^*$ tal que $\sup \operatorname{Re} f(C) \leq \operatorname{Re} f(x_0)$
                \item existe $f \in X^*$ tal que $\sup \operatorname{Re} f(C) < \operatorname{Re} f(x_0)$
                \item no se puede asegurar que $C$ se pueda separar de $x_0$ si el interior de $C$ es vacío
                \item existe un único $f \in X^*$ tal que $\sup \operatorname{Re} f(C) \leq \operatorname{Re} f(x_0)$
            \end{enumerate}
            \item Sea $Z$ un espacio normado de dimensión finita:
            \begin{enumerate}[label={e\arabic*)}]
                \item cualquier aplicación lineal que sale de $Z$ es abierta
                \item cualquier aplicación lineal que llega a $Z$ es continua
                \item cualquier aplicación lineal sobreyectiva que llega a $Z$ es abierta
                \item cualquier aplicación lineal que sale de $Z$ es continua
            \end{enumerate}
        \end{enumerate}
    \end{ejercicio}

    \newpage
    \setcounter{ejercicio}{0} % Reiniciar contador de ejercicios
    \noindent
    \textbf{Solución.}


    \begin{ejercicio}
        En el espacio de Banach $X = (C\left[0, 1\right], \|\cdot\|_\infty)$ se considera el funcional lineal $\varphi$ definido por
        \[
            \varphi(f) = \int_0^1 (2t - 1) f(t) \, dt \quad \forall f \in C\left[0, 1\right].
        \]
        \begin{enumerate}[label=\alph*)]
            \item Pruebe que $\varphi$ es continuo y calcule su norma.
            
            Veamos que $\varphi$ es continuo. Para ello:
            \begin{align*}
                |\varphi(f)| &= \left|\int_0^1 (2t - 1) f(t) \, dt\right| 
                \leq \int_0^1 |2t - 1| |f(t)| \, dt \\
                &\leq \|f\|_\infty \int_0^1 |2t - 1| \, dt
                = \|f\|_\infty \left[\int_0^{\nicefrac{1}{2}} (1 - 2t) \, dt + \int_{\nicefrac{1}{2}}^1 (2t - 1) \, dt\right] \\
                &= \|f\|_\infty \left[\left(t - t^2\right)\Big|_0^{\nicefrac{1}{2}} + \left(t^2 - t\right)\Big|_{\nicefrac{1}{2}}^1\right] = \|f\|_\infty \left[\left(\frac{1}{2} - \frac{1}{4}\right) + \left(1 - 1 - \left(\frac{1}{4} - \frac{1}{2}\right)\right)\right] \\
                &= \|f\|_\infty \left[\frac{1}{4} + \left(-\frac{1}{4} + \frac{1}{2}\right)\right] 
                = \|f\|_\infty \cdot \frac{1}{2}.
            \end{align*}

            Por lo tanto, $\varphi$ es continuo y $\|\varphi\| \leq \nicefrac{1}{2}$. Veamos ahora que $\|\varphi\| = \nicefrac{1}{2}$. La función teórica que deberíamos tomar es $f(t) = \operatorname{sign}(2t - 1)$, pero esta función no es continua. Sin embargo, podemos aproximarla por funciones continuas. Para cada $n \in \bb{N}$, definimos por comodidad:
            \begin{equation*}
                I_n = \left[\frac{1}{2} - \frac{1}{n}, \frac{1}{2} + \frac{1}{n}\right].
            \end{equation*}
            
            definimos la función continua
            \begin{align*}
                f_n(t) = \begin{cases}
                    -1, & t \in [0, \nicefrac{1}{2} - \nicefrac{1}{n}] \\
                    1 & t \in [\nicefrac{1}{2} + \nicefrac{1}{n}, 1] \\
                    \text{lineal} & t \in I_n
                \end{cases}
            \end{align*}

            Notemos que no es necesario dar la expresión explícita de la parte lineal, ya que lo único que nos interesa es que $f_n$ sea continua y que $\|f_n\|_\infty = 1$. Estas funciones se muestran en la Figura~\ref{fig:funciones_fn}. 
            \begin{figure}[h]
                \centering
                \begin{tikzpicture}
                    \begin{axis}[
                        axis lines = middle,
                        xlabel = $t$,
                        ylabel = {$f_n(t)$},
                        ymin = -1.3, ymax = 1.3,
                        xmin = 0, xmax = 1,
                        xtick = {0, 0.5, 1},
                        ytick = {-1, 0, 1},
                        grid = both,
                        width = 11cm,
                        height = 6.5cm,
                        legend pos = south east,
                        legend cell align = left
                    ]

                    \foreach \n/\col in {3/blue, 5/red, 10/teal, 20/orange, 50/purple} {
                        \edef\temp{
                            \noexpand\addplot[color=\col, thick] coordinates {
                                (0, -1)
                                (0.5 - 1/\n, -1)
                                (0.5 + 1/\n, 1)
                                (1, 1)
                            };
                            \noexpand\addlegendentry{$n = \n$}
                        }
                        \temp
                    }

                    \end{axis}
                \end{tikzpicture}
                \caption{Gráfica de la función $f_n$ para distintos valores de $n$.}
                \label{fig:funciones_fn}
            \end{figure}

            Ahora, para cada $n \in \bb{N}$, tenemos que
            \begin{align*}
                \varphi(f_n) &= \int_0^1 (2t - 1) f_n(t) \, dt \\
                &= \int_0^{\nicefrac{1}{2} - \nicefrac{1}{n}} (2t - 1)(-1) \, dt + \int_{I_n} (2t - 1) f_n(t) \, dt + \int_{\nicefrac{1}{2} + \nicefrac{1}{n}}^1 (2t - 1)(1) \, dt.
            \end{align*}

            % // TODO: Continuar
            \item ¿Alcanza $\varphi$ su norma? Justifique su respuesta.
            \item ¿Existe $x^{**} \in X^{**}$ con $\|x^{**}\| = 1$ tal que $x^{**}(\varphi) = \|\varphi\|$? Justifique su respuesta.
        \end{enumerate}
    \end{ejercicio}

    \begin{ejercicio}
        Rodee las afirmaciones verdaderas (que pueden ser una, dos, tres o cuatro):
        \begin{enumerate}[label=\alph*)]
            \item Sea $X$ un subespacio cerrado de $\ell_{3000}$ de dimensión infinita ¿Puede ser $X$ isomorfo a $\ell_1$?
            \begin{enumerate}[label={a\arabic*)}]
                \item Si, porque $\ell_1 \subseteq \ell_{3000}$
                \item No, porque $\ell_1 \nsubseteq \ell_{3000}$
                \item No, porque $\ell_1$ es separable y $X$ no lo es
                \item No, porque $X^*$ es separable y $\ell_1^*$ no lo es
            \end{enumerate}
            \item Sea $\psi \colon \ell_2 \to \bb{R}$ el funcional lineal dado por $\psi(x) = \sum\limits_{n=1}^\infty \frac{x(n)}{n}$ para $x \in \ell_2$:
            \begin{enumerate}[label={b\arabic*)}]
                \item $\psi$ alcanza su norma porque $\ell_2$ es reflexivo
                \item $\psi$ alcanza su norma porque $\ell_2$ es un espacio de Hilbert
                \item $\psi$ no alcanza su norma
                \item $\psi$ no alcanza su norma porque $\ell_2$ no es completo
            \end{enumerate}
            \item Sea $Y = \bb{K}^{(\bb{N})}$ (el espacio vectorial de las sucesiones casi nulas):
            \begin{enumerate}[label={c\arabic*)}]
                \item la norma de la suma y la norma del máximo son equivalentes en $Y$
                \item la norma de la suma y la norma del máximo no son equivalentes en $Y$
                \item cualesquiera dos normas definidas en $Y$ son equivalentes
                \item no se puede definir ninguna norma en $Y$
            \end{enumerate}
            \item En un espacio normado $X$ tomamos un conjunto convexo y cerrado $C$ y un punto $x_0 \notin C$:
            \begin{enumerate}[label={d\arabic*)}]
                \item existe $f \in X^*$ tal que $\sup \operatorname{Re} f(C) \leq \operatorname{Re} f(x_0)$
                \item existe $f \in X^*$ tal que $\sup \operatorname{Re} f(C) < \operatorname{Re} f(x_0)$
                \item no se puede asegurar que $C$ se pueda separar de $x_0$ si el interior de $C$ es vacío
                \item existe un único $f \in X^*$ tal que $\sup \operatorname{Re} f(C) \leq \operatorname{Re} f(x_0)$
            \end{enumerate}
            \item Sea $Z$ un espacio normado de dimensión finita:
            \begin{enumerate}[label={e\arabic*)}]
                \item cualquier aplicación lineal que sale de $Z$ es abierta
                \item cualquier aplicación lineal que llega a $Z$ es continua
                \item cualquier aplicación lineal sobreyectiva que llega a $Z$ es abierta
                \item cualquier aplicación lineal que sale de $Z$ es continua
            \end{enumerate}
        \end{enumerate}
    \end{ejercicio}


\end{document}